\section{Introduction}
\label{ch:introduction}

Since the inception of the Transformer \cite{vaswani2017attention}, large architectures utilizing stacked 
attention mechanisms have dominated the research landscape, exhibiting impressive performance in 
natural language processing tasks. This architectural paradigm scaled rapidly, progressing from hundreds
of millions of parameters to billions, and ultimately to trillions for enterprise-scale models. In parallel, 
with this scaling trend, a significant part of the scientific community focused their efforts on understanding
the internal mechanisms of these models. However, this intense scientific work has focused
almost exclusively towards attention-based architectures, leaving non-transformer based models comparatively less explored. \\ \\
%
Despite their success, Transformers suffer from a fundamental computational bottleneck: the 
self-attention mechanism scales quadratically ($\mathcal{O}(N^2)$) with respect to sequence length, making long-context
processing computationally expensive. This limitation has led to a recent resurgence of interest in recurrent and 
state-based architectures that aim to achieve linear-time efficiency ($\mathcal{O}(N)$) without sacrificing performance. 
Key contributions in this direction include: Structured State Space Models (SSMs) like Mamba \cite{gu2024mambalineartimesequencemodeling}, 
linear transformers such as Receptance Weighted Key Value (RWKV)\cite{peng2023rwkvreinventingrnnstransformer} and 
Extended Long-Short Term Memory(xLSTM) \cite{beck2024xlstmextendedlongshortterm}. \\ \\
%
The xLSTM architecture extends the classic LSTM cell \cite{10.5555/2998981.2999048} by 
introducing exponential gating and matrix-valued memory states with gated updates. These changes look to address the scaling 
problem of the Transformers and the storage bottlenecks of traditional LSTMs, and try to achieve comparable perfromance while 
maintaining linear-time inference.  

\subsection{Problem}

Deep architectures have produced tremendous advancements in the field of machine learning over the past decade. 
Their pattern recognition and information representation capabilities are yet to be completely understood even to this day.
Post-hoc interpretability in deep neural networks is a challenging problem due to several intrinsic mathematical 
and structural properties of these models. \\ \\
%
First, modern architectures operate in high-dimensional representation spaces, 
where features are distributed across multiple dimensions. Moreover, multiple concepts may be encoded within 
overlapping directions of the representation space. \\ \\
%
Second, neural networks apply a sequence of non-linear transformations, making the relationship between input tokens 
and internal representations difficult to describe. As network depth increases, interactions 
between representations become progressively more complex. \\ \\ 
%
Finally, although the computation of hidden states is fully 
specified by the model parameters and architecture, the resulting computation graph is generally too complex to provide 
specific human-interpretable explanations of how particular representations or predictions are produced. \\ \\
%
Beyond interpretability challenges, recurrent-based architectures such as xLSTM also introduce important systems-level 
considerations. In contrast to purely attention-based models, recurrent models maintain an evolving internal state that 
can, in principle, reduce the computational overhead associated with processing long contexts. However, in practice, 
the models suffer from significant generation latency due to the sequential nature of the RNNs. 
This makes efficient inference an important bottleneck to be addressed, particularly when dealing with long sequences or real-time constraints. 

\subsection{Motivation}

The motivation for this work originates from the recent resurgece of recurrent architectures in large-scale sequence 
modelling. While analyizing the current state of research on recurrent neural networks, the xLSTM architecture, which 
Beck et al.\cite{beck2024xlstmextendedlongshortterm} proposed in 2024, became a personal point of interest. As one of the 
first recurrent architectures in recent years to demonstrate competitive performance at scales traditionally 
dominated by Transformer-based models, xLSTM represents an important development in the ongoing search for efficient 
alternatives to self-attention. \\ \\ 
%
While the performance of Transformer-based language models has been extensively studied, a substantial body of 
interpretability research has also emerged to investigate their internal mechanisms. In contrast, the behavior 
of modern recurrent architectures such as xLSTM remains comparatively underexplored. This observation motivated 
the present work, which applies and adapt interpretability techniques inspired by previous studies of 
neural language models. \\ \\
%
By investigating how information is represented, stored, and updated within xLSTM, this thesis aims to contribute 
to a better understanding of the model's memory mechanisms and internal representations. Such insights may help 
validate architectural design choices, identify limitations, and inform future developments of memory-based language models.\\ \\
%
At the same time, it is important to evaluate how alternative architectures perform in practice compared to Transformers, 
which are widely used and benefit from highly optimized inference systems. This motivates the study of xLSTM under realistic 
inference conditions.
\subsection{Contribution}

To address the lack of interpretability analyses for scalable recurrent architectures, the main objective of this thesis is 
to perform a post-hoc evaluation of the internal representations and state dynamics of the xLSTM-7B architecture. 
This work analyzes the internal mechanisms of the model by first studying the hidden states and then examining the individual components
of the xLSTM cells that contribute to the hiden state values.\\ \\ 
%
The first phase of this work establishes an empirical baseline by evaluating the pretrained xLSTM-7B model across standard 
datasets, specifically LAMBADA (OpenAI)\cite{paperno2016lambada}, 
PIQA\cite{bisk2020piqa}, Winogrande\cite{sakaguchi2021winogrande}, ARC-Challenge\cite{clark2018arc}, ARC-Easy\cite{clark2018arc}, 
and Hellaswag\cite{zellers2019hellaswag}. This evaluation represents an assessment of the model's capablities before 
analyzing its internal structure. Following this initial evaluation, the second phase adapts the logit lens\cite{nostalgebraist2020logitlens}
\cite{belrose2023tunedlens} to project intermediate hidden states from the recurrent layers directly onto the vocabulary space. \\ \\
%
Next, we compare how hidden representations evolve across the layers of a recurrent network compared to a Transformer, 
using xLSTM-7B and Mistral-7B. The final two phases of the interpretability study focus on the sequential processing of the architecture. 
Specifically, we examine the evolution of the memory state over time and the contribution of individual 
tokens to long-term memory through the input and forget gates. To analyze these mechanisms, we process
sequences and track the norm of the cell state, together with the activations of the input and forget gates. 
This allows us to observe how information is stored, retained, and updated throughout the sequence. \\ \\ 
%
On top of the interpretability experiments, we evaluate the practical utility of the xLSTM model in a 
question answering setting. We leverage the constant-size memory state of the model and explore caching 
it when processing context that is later reused for answering questions. This approach enables an inference 
speedup that increases with the size of the context, as the stored memory can be reused without recomputing 
the full sequence.
\subsection{Outline}

This thesis is structured as follows. Section~\ref{ch:related-work} presents the theoretical background underlying the work. It introduces Recurrent Neural Networks, Long Short-Term Memory cells, and the Extended Long Short-Term Memory architecture.
\\ \\ 
Section~\ref{ch:method} describes the methodology used in the experiments. More specifically, it outlines the research objectives, the models considered, and the experimental setup, including the theoretical formulation of the conducted experiments. The technologies and implementation details are presented in Section~\ref{ch:implementation}.
\\ \\
Section~\ref{ch:experiments} presents and discusses the experimental results. Finally, conclusions, limitations, and directions for future work are provided in Section~\ref{ch:conclusion-and-future-work}.